% ============================================================================
% PH Foundations v86 -- Section 13 (replaces v85 13.1--13.5) and App. L.6.
% Source: companion repository [2], folder p14d_psion_number_audit/ (PR #11).
% Aligned with the standalone paper paper_baryogenesis/PH_Baryogenesis.tex.
% ============================================================================
%%BEGIN-SEC13
\section*{13  Baryogenesis: Status after the P14 Audit}

\noindent \textbf{Summary. }Earlier releases carried $\eta_B = 6.02\times10^{-10}$ as a derived output at [T2]. The P14 audit withdraws that claim. Equation (1) conserves psion number exactly, field-theory CP is exact, and the Standard-Model sphaleron conversion used in the former eq.~(14) does not apply in the PH. Neither eq.~(1) nor its minimal extension produces the asymmetry (13.3). One candidate mechanism remains, level bifurcation in the pair production of unconfined composites after inflationary reheating (13.7). Its structure follows from symmetry, but its magnitude is not yet explained. This section replaces the former 13.1--13.5.

\subsection*{13.1  What Equation (1) Conserves}

\noindent \textbf{Psion number. }$N_\psi$ is the Noether charge of the overall phase and is exactly conserved by every term of eq.~(1). In the broken phase it is the only exact additive charge (exact charge audit over 20 triples and four condensate variants, zero further candidates).

\noindent \textbf{Continuous internal symmetry. }The full Lie algebra of internal transformations leaving $(\bar\psi J^a\psi)(\bar\psi\Sigma^a\psi)$ invariant is two-dimensional and abelian, $U(1)_N\times U(1)_M$, with $M=\sum m_k$; no continuous chiral extension survives the $G_{\rm UV}$ term. The winding condensate of App.~J.6.2 breaks $U(1)_M$, leaving a residual discrete charge $D=(2M-6Q+3)\bmod 4$.

\noindent \textbf{CP. }The coupling tensor is real and $\varepsilon_\psi$ is CP-even: field-theory CP is exact. The Z6 vacua break CP spontaneously, but conjugate vacua are degenerate ($\Delta C(s)$ symmetric).

\noindent \textbf{Antimatter. }Only one reading is consistent with observed reactions: antiparticles are composites of antipsions ($N_\psi=-3$, charge $-Q$). The index inversion $k\to(6-k)\bmod 6$ (App.~A.4) describes the Z6 label of the antipsion, not a psion composite. The alternative reading (antiparticle $=$ psion triple with inverted indices) fails psion-number conservation already in $\beta^\pm$ decay, fails charge conservation in $n\to pe^-\bar\nu$ and $\mu^-\to e^-\bar\nu_e\nu_\mu$, and gives $\bar\Sigma^+$ the charge $+1$.

\subsection*{13.2  The Sakharov Conditions [16], Revisited}

\noindent \textbf{(1) Baryon-number violation. }Nothing built from psions violates $N_\psi$, including an emergent sphaleron. With $N_\psi/3=B+L$ (13.4), $B+L$ is conserved and only $B-L$ can change: the reverse of the Standard Model. The sphaleron factor 8/23 and the coefficient $C_{\rm core}$ of the former eq.~(14) therefore have no basis in the PH. In addition, $C_{\rm core}=8.4434\times10^{-4}$ could not be reproduced from its stated inputs.

\noindent \textbf{(2) C and CP violation. }Eq.~(1) supplies none (13.1). The former statement that the chirality-flip amplitude $\propto(m_e/\Lambda_{\rm meta})^2$ is a CP-odd phase is withdrawn: a real amplitude carries no phase.

\noindent \textbf{(3) Departure from equilibrium. }Available: the formation of composites is out of equilibrium when the contact rate is comparable to the expansion rate (13.7). Temperature histories alone redistribute abundances among cascade levels but never create an asymmetry.

\subsection*{13.3  Withdrawn Results and Closed Mechanisms}

\noindent \textbf{Withdrawn. }(i) The former eqs.~(14)--(15), $\eta_B = C_{\rm core}(g_{1/2}/g_S)\varepsilon_\psi = 6.02\times10^{-10}$: the sphaleron input has no PH basis, and $C_{\rm core}$ is not reproducible. (ii) Eqs.~(37)--(38) of App.~J: $\xi^*=434.09$ was computed from the target $\eta_B$, so eq.~(38) is an identity, not a check. Independent routes to $\xi^*$ give $\xi^*\le1.08$ (Potts Z3; second-order NJL transition with $T_c=9.6$ MeV) and $\xi^*=3.6$--$39$ (Coleman--Weinberg with thermal inflation). (iii) The claim that the PH satisfies all three Sakharov conditions without additional ingredients. (iv) Prediction F12 is suspended (§9).

\noindent \textbf{Reclassified. }$\varepsilon_\psi=25/169$ remains a result of SU(2)$_{\rm meta}$ geometry (App.~J, conditional on P23). Its link to $\eta_B$ is not established.

\noindent \textbf{Closed mechanisms. }Each was tested under a pre-registered definition, and each fails for a stated reason:
\begin{itemize}
\item Emergent sphalerons: they conserve $N_\psi$.
\item Z6 domains with spontaneous CP: they form at $T\approx3$ MeV, when assembly is frozen ($\Gamma/H=3\times10^{-37}$). Their maximal comoving size today is $\sim30$ pc, against the $\gtrsim1$ Gpc required of antimatter regions [19]. Degenerate conjugate vacua give a zero global average.
\item Temperature histories: no asymmetry from symmetric initial data.
\item CP terms built from PH operators: 16 commutator and rearrangement patterns, alone or in pairs, fit the observed composition no better than random directions.
\item A phase inserted in $J\!\cdot\!\Sigma$: it generates an asymmetry, but exotic triples receive 2--2500 times the proton asymmetry. Without a conserved label, no CP source selects the observed composites.
\item Species mixing in the minimal extension of 13.4: the proton lifetime bounds its coefficient to $c\le1.3\times10^{-8}$, which places the active window above $2.7\times10^{15}$ GeV, outside the contact description.
\end{itemize}

\subsection*{13.4  Baryon Number and the Level Rule}

\noindent \textbf{What observation requires. }With the cascade level treated as a quantum number, the observed reactions $\beta^\pm$ decay, $e^-$ capture, $\pi^+\to\mu^+\nu_\mu$, $\mu$ decay, $\nu_\mu n\to\mu^-p$, $\bar\nu_ep\to e^+n$, $\Lambda\to p\pi^-,n\pi^0$, $\Sigma^0\to\Lambda\gamma$ and neutrino oscillations require exactly two additive charges: $B$ (states of the level-8 class) and $L$ (level-10 and neutrino classes), with $N_\psi/3=B+L$.

\noindent \textbf{What eq.~(1) allows. }With only $N_\psi$ and $Q$ conserved, seven unobserved processes are allowed: $p\to\nu\pi^+$, $n\to\nu\pi^0$, $n\to3\nu$, $p\,e^-\to\nu_\mu\nu_\tau$, and three level jumps within one triple ($n\to\nu_\tau\gamma$, $\Lambda^0\to\nu_e\gamma$, $\Sigma^0\to\nu_\mu\gamma$). All have $\Delta B=-\Delta L$. Against the measured limits ($p\to\bar\nu\pi^+>3.9\times10^{32}$ yr and $n\to\bar\nu\pi^0>1.1\times10^{33}$ yr [17]; $n\to$ invisible $>5.8\times10^{29}$ yr [18]), the required amplitude suppression is $\lesssim10^{-8}$ relative to a contact operator at $\Lambda_{\rm meta}$. For same-triple jumps it is $\lesssim5\times10^{-30}$ relative to an electromagnetic transition (reference lifetime $10^{29}$ yr).

\noindent \textbf{What eq.~(1) cannot supply. }$B$ is not a Noether charge: it lies outside the span of $(M,Q,N_\psi)$. It is not topological: the vacuum manifold has $\pi_2=\pi_3=0$, so no skyrmion-like baryon number exists [21]. And the level is not an internal degree of freedom, since states are labelled by the triple alone. Hence the dual-sector identification of \S6 and L.4.4, in which the same triple is either a baryon or a neutrino ((0,1,5) $=n/\nu_\tau$, (0,2,4) $=\Sigma^0/\nu_\mu$, (1,2,3) $=\Lambda^0/\nu_e$), is incompatible with neutron stability as eq.~(1) stands. Correction to \S6: there are five neutral three-body triples, not three. They are (0,1,4), (0,1,5), (0,2,4), (0,2,5) and (1,2,3); the two additional ones carry $D=2$.

\noindent \textbf{Minimal extension (registered, not adopted). }A two-valued species index, $\psi\to\psi_s$ with $s\in\{B,L\}$ and the currents of eq.~(1) summed over species, preserves $U(1)_B\times U(1)_L$ if $SU(2)_{\rm species}$ is broken. All observed reactions are allowed, and all seven unobserved processes are forbidden exactly. Its costs:
\begin{itemize}
\item The $L$ sector contains six positively charged matter leptons that must be heavy.
\item The numerical predictions survive only if the species index does not enter $|\beta_0|=N_{\rm Dirac}+2/N_f$. Doubling $N_f$ gives $m_e=0.065$ MeV.
\item A label carried by the constituents cannot be chosen at formation.
\end{itemize}
This is recorded as open problem \textbf{P14-B}: a conserved baryon (rung) label.

\subsection*{13.5  Link to $N_{\rm eff}$ (status note)}

\begingroup\centering $\Delta$Neff\_PH = 3 $\cdot$ (15/7$\pi$\textsuperscript{2}) $\cdot$ ($\mu$$\nu$/T)\textsuperscript{2} = 0.010,   Neff\_PH = 3.054.\hspace{2em}(16)\par\endgroup

\noindent Equation (16) is retained, but its shared-parameter status with $\eta_B$ lapses with 13.3. If $\mu_\nu/T=0.123$ is read as a lepton-number asymmetry, it conflicts with psion-number bookkeeping. With total $N_\psi=0$, $B=-L$ and charge neutrality, the antineutrino excess is two per hydrogen atom, a degeneracy $\xi_\nu\approx2.5\times10^{-9}$, some $5\times10^7$ times smaller than 0.123. Equation (16) must therefore rest on a different reading of $\mu_\nu$ or be reclassified.

\subsection*{13.6  Structural Predictions}

\noindent \textbf{(a) }No nucleon decay with $\Delta B=\Delta L$: $p\to e^+\pi^0$ and similar modes change $N_\psi$ and are forbidden exactly [T1]. \textbf{(b) }Any nucleon decay, if observed, has $\Delta B=-\Delta L$ ($p\to\nu\pi^+$, $n\to\nu\pi^0$, $n\to3\nu$). \textbf{(c) }No $0\nu\beta\beta$ and no $n$--$\bar n$ oscillation, since both change $N_\psi$. \textbf{(d) }If the asymmetry formed with total $N_\psi=0$: a cosmic antineutrino excess of two per hydrogen atom. \textbf{(e) }Dark matter as a thermal relic of the cascade top would require level $\approx4.8$ ($\approx56$ TeV); the canonical $m_{\rm DM}=7.2\times10^{10}$ GeV requires a non-thermal origin (§14.3).

\subsection*{13.7  A Candidate Mechanism: Level Bifurcation at Reheating}

\noindent \textbf{Unconfined spectrum. }By the confinement rule of App.~H/I (summed index $I \bmod 6\in\{0,3\}$), exactly eight triples are unconfined: (0,1,2), (0,1,5), (0,2,4), (0,4,5), (1,2,3), (1,3,5), (2,3,4) and (3,4,5). These are the triples with one constituent of each proto-colour, and they are the states identified with observed particles. All fractional-charge triples and both $D=2$ neutrals are confined. Exotic triples therefore do not constrain any mechanism below $\Lambda_{\rm meta}$.

\noindent \textbf{Postulates. }Below $\Lambda_{\rm meta}$ only these composites and radiation exist, so composites are pair-produced from radiation. Each triple has four channels, and $t_B\bar t_L$ and $t_L\bar t_B$ are CP conjugates from a self-conjugate initial state, the minimal setting for CP violation compatible with CPT and unitarity [20]. (P1) The rung (baryonic $t_B$ or leptonic $t_L$) is fixed at production and conserved by every process except annihilation back into radiation. (P2) The conjugate channels are biased by $\varepsilon_t$, with source $\varepsilon_t$ times the net production rate, so the source vanishes in equilibrium. (P3) Production occurs after reheating at $T_{\rm RH}<\Lambda_{\rm meta}$ with the natural coupling, so a single condensate vacuum covers the observable Universe.

\noindent \textbf{Structural consequences ([T1] given P1--P2). }Each production event changes $B$ and $L$ oppositely and creates a triple with its antitriple. Hence $B=-L$, total charge zero and $\Delta N_\psi=0$, and the late-time state is hydrogen plus two net antineutrinos per baryon.

\noindent \textbf{CP source. }The vacuum phase of the winding condensate changes $M$ by two units, and $U(1)_M$ is exact in eq.~(1); since the condensate itself generates the bond of J.6.3, a $U(1)_M$ rotation removes the phase exactly in every vacuum. The phase is physical only relative to an explicit breaking of strength $\delta$, with CP-odd effects $\propto\delta^2$. The only explicit breaking in the PH, $\kappa_{\rm CP}=(m_e/\Lambda_{\rm meta})^2$ (K.10), gives a bias $\sim10^{-62}$. The PH as written has no CP source of adequate size.

\noindent \textbf{Efficiency and open problem. }With the natural coupling, $\eta_B/\varepsilon$ reaches $\sim0.1$ for $T_{\rm RH}\sim4$--$7\times10^9$ GeV, and the asymmetry is washed out above $\sim1.5\times10^{10}$ GeV. The result is insensitive to gradual reheating (peak temperature up to 60 $T_{\rm RH}$). The observed $\eta_B$ needs a bias of $10^{-7}$--$10^{-6}$, far above what the PH supplies, and the sign depends on the binding law and on $T_{\rm RH}$. The reheating temperature is not a prediction.

\noindent \textbf{Status ([T3]) and conditions. }(C1) Derive the rung label, for instance as a winding charge changed only through the deconfined core, where the action per mode $S/N_e$ gives a suppression $e^{-21\Delta\ell}$, enough for 13.4. (C2) Identify a derived explicit breaking of $U(1)_M$ (or another CP-odd ingredient) of strength $\delta\sim10^{-3}$. (C3) Derive the pair-production rate of composites.
%%END-SEC13
%%BEGIN-L6
\subsection*{L.6  v86 Revision Record: the P14 Audit (October 2026)}

\noindent All work packages were pre-registered, with SHA-256 and a UTC timestamp, before execution. Failures and corrections are recorded with the same weight as results; the full log is in \path{p14d_psion_number_audit/README.md} of repository [2].

\noindent \textbf{L.6.1 Changes. }§13 rewritten (13.1--13.7); $\eta_B$ derivation withdrawn; F12 suspended; App.~J.1 and J.4 annotated; App.~A.4 antimatter statement corrected (antipsions); §8.2 status note: the missing level rule is a conserved baryon label (P14-B); §14.3 direct route marked unsupported; open-problem tables gain P14-B; abstract, §4.1, §12 and §15--19 summaries updated accordingly. No other numerical prediction changes.

\noindent \textbf{L.6.2 Errors made during the audit and corrected. }(i) Sphalerons were first treated as external to the PH; they are psion composites and conserve $N_\psi$. (ii) Antimatter was first modelled as antipsions, then switched to index-inverted psion triples on a reading of A.5, and restored to antipsions by the reaction test of 13.1; the second step was the wrong one. (iii) In the domain test, unitarity was first imposed jointly across processes with different density scaling, which created asymmetry in equilibrium. A control caught this, the test was corrected, and the faulty version is retained. (iv) In the calibration test, the dark-matter block first patched a copy of the module globals and was redone in an addendum. (v) The anticipated count of neutral triples was wrong (it is five). (vi) In the bifurcation test, the registered endpoint bookkeeping had a sign error, which was corrected (two antineutrinos per baryon), and the anticipated efficiency was wrong by about three orders of magnitude. (vii) The bifurcation networks first used free psions below $\Lambda_{\rm meta}$, which App.~I forbids; the mechanism was reformulated as pair production of unconfined composites. (viii) Rung assignments in two tests contradicted the stability requirement for charged leptonic copies; corrected, they make the sign of $\eta_B$ law-dependent. (ix) A test of the vacuum phase counted the condensate-generated bond twice, as an independent real term; without the double counting the phase is removable by $U(1)_M$, and the claimed CP source does not exist.

\noindent \textbf{L.6.3 Registered control failures. }All are at or near the floating-point floor or due to integrator limits: assembly-network controls ($2\times10^{-13}$ against $10^{-15}$); temperature-programme step convergence (1--2\%); domain test C3 ($2.7\times10^{-6}$) and C4 ($1.3\times10^{-12}$); phase-in-$J\!\cdot\!\Sigma$ test C3 ($1.8\times10^{-6}$) and one sensitivity run not completed; bifurcation tests, $B+L=0$ at the numerical floor ($\le1.1\times10^{-8}$).
%%END-L6
